\section{An axiom set that was not minimal}

The interpretation was first stated with four frame conditions. A check that
drops each condition and re-runs the proofs reported two of them inert: dropping
either lost no theorem --- where ``lost'' then meant ``no longer proved'', which is
the wrong criterion and is corrected below. Neither was a reprieve. Dropping \emph{both} loses three
theorems, because each was derivable from the other three, so the four were a
redundant presentation of the same constraint. The interpretation now rests on
three independent conditions --- every coordinate supports the certificate, every
edge runs from a coordinate to the certificate, and the certificate is not one
of its own coordinates --- and the fourth, that the certificate supports nothing,
is discharged as a theorem about them. Each of the three loses at least one
theorem when dropped, and which theorem each carries is recorded.

An axiom that no theorem needs is either decoration or a redundant presentation
of an axiom that is doing the work, and the two are indistinguishable from
inside. The check that told them apart is worth more here than the axiom it
removed.

That check has since been wrong twice, and both corrections are worth stating
because they change what ``load-bearing'' can mean rather than only how it is
computed. It first counted any theorem that stopped being \emph{proved} when an axiom
was dropped. That credits an axiom for the solver failing to settle a question:
a solver returning \texttt{unknown} has said nothing about whether the theorem still
holds. Requiring an actual countermodel fixed the criterion and broke the
stability, because refuting a universally quantified claim over an uninterpreted
sort is an unbounded model search --- the same dropped axiom yielded a
countermodel on one run and \texttt{unknown} on the next. The measurement now asks for
a countermodel in a universe bounded to four nodes, which is sound in that
direction and in no other (a countermodel in a small universe is a countermodel;
failing to find one there proves nothing), and takes the intersection over
repeated runs.

Under that strictest reading all three conditions remain necessary, so the
conclusion survived both corrections. The per-condition detail did not: the
edge-restriction condition refutes between one and three theorems on identical
repeated runs, of which exactly one falls every time, and only that one is
credited to it.
